% ----------------------------------------------------------
% Glossário
% ----------------------------------------------------------

%Consulte o manual da classe abntex2 para orientações sobre o glossário.

%\glossary




% ----------------------------------------------------------
% Glossário (Formatado Manualmente)
% ----------------------------------------------------------

\chapter*{GLOSSÁRIO}
\addcontentsline{toc}{chapter}{GLOSSÁRIO}

{ \setlength{\parindent}{0pt} % ambiente sem indentação

\textbf{Anycast}: Técnica de rede onde a mesma faixa de endereços IP é utilizada em múltiplos servidores em locais diferentes. As requisições são roteadas para o servidor topologicamente mais próximo do cliente, melhorando a performance e a resiliência.

\textbf{Broker}: Em sistemas distribuídos, um intermediário que gerencia a comunicação entre diferentes componentes de software. Ele recebe mensagens de um publicador e as direciona para os assinantes interessados.

\textbf{Buffer}: Uma área de memória temporária usada para armazenar dados enquanto eles estão sendo movidos de um lugar para outro, compensando diferenças de velocidade entre o produtor e o consumidor dos dados.

\textbf{Cache}: Componente de hardware ou software que armazena dados para que futuras requisições por esses mesmos dados possam ser atendidas de forma mais rápida; os dados armazenados em um cache podem ser o resultado de uma computação anterior ou uma cópia de dados armazenados em outro lugar.

\textbf{Cloud}: Modelo de computação que permite o acesso sob demanda, via rede, a um conjunto compartilhado de recursos computacionais configuráveis (ex: redes, servidores, armazenamento, aplicações e serviços).

\textbf{Contêiner}: Uma unidade padrão de software que empacota o código e todas as suas dependências para que a aplicação seja executada de forma rápida e confiável de um ambiente computacional para outro.

\textbf{Cookies}: Pequenos arquivos de texto que os sites enviam para o navegador de um usuário para armazenar informações sobre sua sessão, preferências ou estado de autenticação, permitindo uma experiência de navegação personalizada.

\textbf{Datacenter}: Instalação física centralizada que abriga a infraestrutura de TI de uma organização, incluindo servidores, sistemas de armazenamento de dados, equipamentos de rede e outros componentes necessários para operar e gerenciar aplicações e dados.

\textbf{Docker}: Plataforma de software de código aberto usada para desenvolver, empacotar e executar aplicações em contêineres, permitindo a separação de suas aplicações da infraestrutura para que você possa entregar software rapidamente.

\textbf{Ethernet}: Família de tecnologias de rede de computadores com fio para redes locais (LANs), definindo o cabeamento e os sinais elétricos para a camada física, e um formato de pacote e protocolo para a camada de enlace de dados do modelo OSI.

\textbf{Failover}: Processo de comutação automática para um sistema ou componente de backup em caso de falha do sistema primário, garantindo a continuidade do serviço com o mínimo de interrupção.

\textbf{Firewall}: Dispositivo de segurança de rede que monitora o tráfego de rede de entrada e saída e decide permitir ou bloquear tráfegos específicos com base em um conjunto definido de regras de segurança.

\textbf{Gateway}: Nó de rede que conecta duas redes diferentes que usam protocolos de comunicação distintos, atuando como um ponto de entrada e saída para o tráfego de rede.

\textbf{Host}: Qualquer dispositivo computacional conectado a uma rede que oferece recursos, serviços e aplicações aos usuários ou a outros nós na rede.

\textbf{Hypervisor}: Software, firmware ou hardware que cria e executa máquinas virtuais (VMs). Ele permite que um único computador host suporte múltiplas VMs, compartilhando virtualmente seus recursos, como processador e memória.

\textbf{Middleware}: Software que atua como uma ponte entre um sistema operacional e as aplicações nele executadas, oferecendo serviços comuns e funcionalidades de conectividade que não são fornecidas nativamente pelo sistema operacional.

\textbf{Mininet}: Emulador de rede que cria uma rede de hosts virtuais, switches, controladores e links em um único kernel Linux, permitindo o desenvolvimento, teste e prototipagem de aplicações de Redes Definidas por Software (SDN).

\textbf{Namespaces}: Recurso do kernel Linux que particiona recursos do kernel de forma que um conjunto de processos veja um conjunto de recursos, enquanto outro conjunto de processos vê um conjunto diferente, sendo a base tecnológica dos contêineres.

\textbf{NetBox}: Aplicação web de código aberto projetada para gerenciamento de infraestrutura de rede (IPAM - IP Address Management) e de infraestrutura de datacenter (DCIM - Data Center Infrastructure Management).

\textbf{OpenFlow}: Protocolo de comunicação que permite que o plano de controle de uma rede definida por software (SDN) interaja e gerencie os dispositivos do plano de dados (como switches e roteadores).

\textbf{Peering}: Interconexão voluntária de redes de Internet administrativamente separadas com o propósito de trocar tráfego entre os usuários de cada rede.

\textbf{Player (aplicação Player MPEG-DASH)}: Aplicação cliente responsável por requisitar os segmentos de mídia de um servidor HTTP e decodificá-los para reprodução, adaptando a qualidade do vídeo com base nas condições da rede.

\textbf{Proxy}: Servidor que atua como intermediário para requisições de clientes buscando recursos de outros servidores. Ele pode ser usado para filtrar requisições, registrar logs, ou compartilhar conexões.

\textbf{Python}: Linguagem de programação de alto nível, interpretada, de script, imperativa, orientada a objetos, funcional, de tipagem dinâmica e forte, conhecida por sua sintaxe clara e legibilidade.

\textbf{Remote}: Refere-se a um sistema, dispositivo ou recurso que está localizado em um local diferente do usuário ou do sistema principal, sendo acessado através de uma rede.

\textbf{Round-Robin}: Algoritmo de escalonamento de processos ou de balanceamento de carga que distribui as tarefas ou o tráfego de rede de forma sequencial e cíclica entre os recursos disponíveis.

\textbf{Rust}: Linguagem de programação multiparadigma compilada, focada em segurança, especialmente a segurança de concorrência e de memória, mantendo um alto desempenho.

\textbf{Script}: Sequência de comandos armazenada em um arquivo de texto simples, executada por um interpretador, para automatizar tarefas ou configurar um sistema.

\textbf{Shell}: Programa de interface de usuário, geralmente baseado em linha de comando, que permite aos usuários controlar o computador através da inserção de comandos textuais.

\textbf{Socket}: Ponto final de um fluxo de comunicação bidirecional entre processos em uma rede. É definido por um endereço IP e um número de porta.

\textbf{Stalls}: Em streaming de vídeo, uma interrupção na reprodução causada quando o buffer do player fica vazio, forçando o vídeo a pausar enquanto mais dados são baixados.

\textbf{Startup}: O processo de inicialização de um sistema computacional ou de uma aplicação, envolvendo o carregamento do sistema operacional, serviços e programas necessários para sua operação.

\textbf{Streaming}: Tecnologia de transmissão contínua de dados, geralmente áudio e vídeo, de um servidor para um cliente, permitindo que a reprodução comece antes que todo o arquivo tenha sido baixado.

\textbf{Switches}: Dispositivos de rede que operam na camada de enlace de dados (camada 2 do modelo OSI) para conectar múltiplos dispositivos em uma mesma rede local (LAN), encaminhando pacotes de dados apenas para o destinatário específico.

\textbf{WebSockets}: Protocolo de comunicação que proporciona canais de comunicação full-duplex sobre uma única conexão TCP, permitindo a interação em tempo real entre um navegador e um servidor web.

\textbf{Windows}: Família de sistemas operacionais gráficos desenvolvidos, comercializados e vendidos pela Microsoft, amplamente utilizados em computadores pessoais.
}


